% Source: https://artofproblemsolving.com/community/c3h1367344_scalene_quadralateral
% Stable identity: aops-topic-1367344-q01
\textbf{Varignon's theorem.} Let $ABCD$ be the quadrilateral and $K,L,M,P$ the midpoints of $AB,BC,CD,DA$. By the midsegment theorem in triangles $ABC$ and $ACD$, $KL$ and $PM$ are both parallel to $AC$ with length $\dfrac{1}{2}AC$; similarly $LM$ and $KP$ are parallel to $BD$ with length $\dfrac{1}{2}BD$. So $KLMP$ is a parallelogram with sides $\dfrac{1}{2}AC$ and $\dfrac{1}{2}BD$.

Hence $KLMP$ is equilateral (a rhombus) \emph{exactly when the diagonals are equal}, $AC=BD$. This does not force equal sides, so the answer is \textbf{sometimes}.

\textbf{Example.} Take $A=(-2,0)$, $B=(0,-3)$, $C=(6,0)$, $D=(0,5)$. The diagonals $AC$ and $BD$ both have length $8$, and the sides are
\[AB=\sqrt{13},\quad BC=\sqrt{45},\quad CD=\sqrt{61},\quad DA=\sqrt{29},\]
all different, so $ABCD$ is scalene (and convex, since the diagonals cross inside it).

\begin{asy}
usepackage("amsmath");
size(5.5cm);
defaultpen(fontsize(10pt));
pair A=(-2,0),B=(0,-3),C=(6,0),D=(0,5);
pair K=(A+B)/2,L=(B+C)/2,M=(C+D)/2,P=(D+A)/2;
draw(A--B--C--D--cycle,linewidth(0.9));
draw(A--C,gray+dashed); draw(B--D,gray+dashed);
draw(K--L--M--P--cycle,red+linewidth(0.9));
dot(K,red); dot(L,red); dot(M,red); dot(P,red);
label("$A$",A,W); label("$B$",B,S); label("$C$",C,E); label("$D$",D,N);
label("$K$",K,SW,red); label("$L$",L,SE,red); label("$M$",M,NE,red); label("$P$",P,NW,red);
label("\mbox{}",(min(currentpicture,user=true)+max(currentpicture,user=true))/2,black); // keeps the text after the figure black when LaTeX typesets the labels inline
\end{asy}

Its midpoint quadrilateral $KLMP$ has all four sides equal to $4$, so it is a rhombus. It is not always equilateral: for a scalene quadrilateral with $AC\ne BD$ (e.g.\ move $D$ to $(0,4)$, giving sides $\sqrt{13},\sqrt{45},\sqrt{52},\sqrt{20}$ and diagonals $8\ne7$) the sides of $KLMP$ are $4$ and $\dfrac{7}{2}$.
